\section{项目亮点}

\subsection{亮点一：围绕学习动作组织数据，而不是堆叠模型功能}

作品将识别后的内容组织为卡片，是因为后续多个动作需要同一份可寻址材料。卡片既保留适合生成提示词的摘要，也保留用于查阅的正文；既有标签和主公式，也有来源编号。命题结果记录使用了哪些卡片，追问记录命中了哪些卡片，人工修订保留字段级历史。由此，生成、检索和复习不再是相互独立的演示按钮。

这种组织方式解决了两个常见失败。第一，模型回答无法核对来源。若只把整页 OCR 文本临时塞入提示词，用户看到答案后无法回到具体知识点。第二，识别错误会沿后续链路放大。若不能编辑卡片，一次公式误读会持续影响出题和追问。当前实现允许先修正卡片再生成，并把修改前后的值留档。它不能提高底层视觉模型本身的准确率，但能降低错误继续传播的成本。

与“上传后一次性生成整套课程”的方案相比，卡片化增加了用户确认步骤，也增加数据库和编辑界面的复杂度。这个取舍是有意的。对于学习材料，可修改和可追溯比一次生成更多文字更重要。系统没有自动替用户决定所有知识分类，不确定内容保留“无法识别”，用户可以在进入练习前纠正。

由此，来源、卡片、题目、作答和事件分别承担后续消费与错误定位，具体数据关系见第 4 节的数据模型图。

\subsection{亮点二：生成结果拥有独立的否定路径}

语言模型可以在同一段上下文中同时给出题和“已检查”的声明，但这种声明缺少独立性。本作品要求命题模型提交可计算账单，再由不读取其推理过程的 sympy 代码核算。质检只看到表达式、代入值和期望值。它的作用不是理解教学语言，而是判断数值等式是否成立。对于能够表示成这类等式的题，错误可以被抓到并返回生成阶段。

这种设计的重点不是多调用一个角色，而是把否定权交给不同能力。如果质检仍由同一个语言模型完成，它可能复制命题时的错误。如果只做字符串比较，等价表达式又容易被误判。符号计算提供了与文本生成不同的失败模式，两者组合比单一路径更有判别力。

\begin{figure}[H]
\centering
\resizebox{0.97\textwidth}{!}{%
\begin{tikzpicture}[
  node distance=8mm and 11mm,
  state/.style={draw=ink!60,fill=palegreen,rounded corners=3pt,align=center,font=\small,minimum width=2.55cm,minimum height=1.05cm,inner sep=5pt},
  decision/.style={draw=amber!75,fill=amber!8,diamond,aspect=2.2,align=center,font=\small,inner sep=2pt},
  terminal/.style={draw=ink,fill=ink,text=white,rounded corners=3pt,align=center,font=\small,minimum width=2.4cm,minimum height=1.0cm},
  arr/.style={-{Stealth[length=2.2mm]},draw=muted,thick},
  lbl/.style={font=\scriptsize,text=muted,fill=white,inner sep=1.5pt}
]
\node[state] (input) {卡片与练习要求};
\node[state,right=of input] (compose) {命题\\第 $r$ 轮};
\node[state,right=of compose] (verify) {质检\\逐题生成账单};
\node[decision,right=of verify] (pass) {全部通过？};
\node[terminal,right=of pass] (deliver) {按级别交付\\写入题库};
\node[state,below=1.45cm of pass] (drop) {筛除未通过项\\必要时结构降级};
\node[terminal,below=1.45cm of compose] (reject) {无合格题\\拒绝入库};
\draw[arr] (input) -- (compose);
\draw[arr] (compose) -- node[lbl,above]{questions + verify} (verify);
\draw[arr] (verify) -- (pass);
\draw[arr] (pass) -- node[lbl,above]{是} (deliver);
\draw[arr] (pass.south) -- node[lbl,right]{$r<3$} (drop.north);
\draw[arr] (drop.west) -| node[lbl,pos=0.25,above]{带原因与轮次重生成} (compose.south);
\draw[arr] (drop.east) -| node[lbl,pos=0.28,above]{$r=3$ 且无合格题} (reject.east);
\draw[arr] (drop.north east) to[bend left=12] node[lbl,right]{保留通过当前检查者} (deliver.south);
\end{tikzpicture}}
\caption{命题与质检的有限反馈状态图。通过、驳回、重试、筛除和降级均产生真实事件；达到上限后不允许无限循环。}
\label{fig:verify-loop}
\end{figure}

有限重试是另一个关键约束。坏提示词或不稳定的模型可能持续生成不合规输出，无限重试会消耗额度并占住请求。系统将 \env{VERIFY_MAX_RETRY=2} 解释为首轮之后最多两次重试。轮次耗尽时，优先保留已通过当前检查的题；v0.1.6 在没有任何通过项时允许把结构合格题标为“结构检查”后交付，并通过显式事件说明发生了降级。关闭 \env{VERIFY_FALLBACK_ON_EXHAUST} 可恢复“全不合格则拒绝”的旧行为。无论采用哪条路径，界面都不能把结构检查写成符号验算。

\subsection{亮点三：正确性声明分级，不用一个徽章覆盖所有情况}

生成题的可信程度并不相同。带合法账单的题可以完成数值代入，无账单题只能检查结构，关闭结构检查时的旧路径则没有有效检查。系统把三种情况分别记为 \env{symbolic}、\env{structural}、\env{unchecked}，在数据库、结束事件和界面徽章中保持一致。这样，用户不必从“通过”两个字猜测系统到底做了什么。

分级还改变了失败处理。真实模型有时把 \env{a*c} 写成 \env{ac}，或在表达式里残留 LaTeX。这些更像格式能力不足，而不是答案必然错误。当前版本允许将有限类别的格式问题降级为结构检查，但注入字符、非数值代入、危险幂结构等安全边界不参与这种宽松处理。降级原因进入日志，不能悄悄消失。

\begin{table}[H]
\centering\footnotesize
\caption{三级检查语义及其能够支持的结论}
\begin{tabularx}{\textwidth}{@{}L{2.5cm}L{3.7cm}Y L{2.7cm}@{}}
\toprule
\textbf{级别} & \textbf{执行内容} & \textbf{可以说明什么} & \textbf{界面含义} \\
\midrule
\env{symbolic} & 受限解析、数值代入、容差比较 & 提交的等式在给定代入下自洽；不证明题意没有歧义，也不等于所有数学命题已证明 & 已验算 \\
\env{structural} & 检查题干、答案、题型，选择题另查选项与答案范围 & 对象形状完整，可进入界面；不保证答案内容正确 & 结构检查 \\
\env{unchecked} & 关闭结构检查后的旧行为 & 没有形成足够的质量证据，只保留兼容路径 & 未验算 \\
\bottomrule
\end{tabularx}
\end{table}

\subsection{亮点四：过程可见，错误可以定位与恢复}

七个 Agent 的状态通过统一事件协议进入看板，命题轮次、驳回对象和验算账单均可查看。业务驳回使用 \env{rejected}，系统异常使用 \env{error}，颜色和符号同时区分，避免把正常质检误认为崩溃。

HTTP 操作立即返回 \env{run_id}，服务端先建运行再启动协程，迟到订阅者从历史起点补事件。断流后的重试依次读取内存历史、落库事件，都无记录才重新 POST，减少重复上传和作答。数据库写失败不阻塞当前 SSE；单 run 内存上限 500 条，落库保留最近 200 个运行，以有限历史换取可预测占用。

\subsection{亮点五：无密钥复现不是另一套假界面}

评审环境可能存在网络、账号或时间限制。缺省 mock 让上传、命题驳回、判分讲解、检索出处和批改都可以重复出现。mock 数据固定首轮第 4 题的期望值错误，使驳回在每次演示中稳定发生；第二轮恢复正确值。这个样例不是预先写好的“驳回动画”，实际判定仍由验算函数计算，流水线测试还断言命题恰好执行两次。

固定样例提高了可复现性，但不能证明真实模型表现。作品在界面、README 和本文中同时说明 mock 不读图，真实 Provider 才进行视觉识别。测试也强制 mock 与临时数据库，防止开发机环境变量里的真实 Key 让离线测试意外访问网络。这样的测试设计让“全绿”具有明确含义：业务规则通过，而不是外部服务当时可用。

\subsection{亮点六：轻量交付与可选增强保持分离}

缺省检索使用本地 2-gram，没有向量数据库和嵌入服务依赖。需要改善同义改写召回时，可显式启用本地 BGE-M3 与 RRF 融合，端点故障自动退回基线。缺省 Web 模式便于开发，桌面目录版便于演示，便携版便于单文件分发。三种形态共享业务代码，差异集中在启动和配置定位。

这一原则也体现在没有引入的依赖上。编排只有一个反馈回路，因此使用 \env{asyncio} 状态机而不引入 LangGraph；个人知识库规模下用 SQLite 和 Python 余弦计算，不引入原生向量扩展；函数图采用自写表达式解析和 SVG，不引入大型绘图库。少依赖不是目标本身，判断依据是新增依赖能否解决当前已经发生的问题。

\begin{table}[H]
\centering\footnotesize
\caption{主要设计取舍及代价}
\begin{tabularx}{\textwidth}{@{}L{2.8cm}L{4.3cm}Y@{}}
\toprule
\textbf{决策点} & \textbf{当前选择} & \textbf{收益、代价与放弃方案的原因} \\
\midrule
编排框架 & 自研 \env{asyncio} 编排器 & 状态少、代码可直接定位；复杂分支扩展能力弱。赛程内引入框架的学习与调试成本高于收益 \\
实时协议 & SSE 而非 WebSocket & 服务端单向推送足够，协议简单；不自动重连，通过历史与落库回放恢复 \\
生成质检 & 本地 sympy 与结构检查 & 获得独立否定路径；只覆盖可表示为数值账单的题，开放题仍需其他机制 \\
检索基线 & 2-gram 字段加权 & 零依赖、可离线；同义与跨语能力有限。向量只作为显式增强，故障可回退 \\
图形渲染 & 自写解析器与 SVG & 无 \env{eval}、无大型绘图库；支持的数学语法和模板数量有限 \\
数据存储 & 单进程 SQLite + WAL & 部署简单、来源关系清楚；不做高并发声明，多 worker 的内存看板也不一致 \\
桌面分发 & 目录版优先，便携版保留 & 免除环境安装；体积大且未签名，便携版自解压影响首次体验 \\
\bottomrule
\end{tabularx}
\end{table}

\subsection{与常见工具路径的差异}

2026 年 11 月 15 日的 GitHub 公开仓库检索中，“错题本 ai”命中 237 个，“quiz generation answer verification”命中 6 个。该检索不含私有项目和商业产品，只用于确定差异点：错题本并不稀缺，答案的独立验证值得单独实现。

\begin{table}[H]
\centering\footnotesize
\caption{常见工具路径与本作品的职责差异}
\begin{tabularx}{\textwidth}{@{}L{2.7cm}L{3.7cm}L{3.7cm}Y@{}}
\toprule
\textbf{维度} & \textbf{公式识别服务} & \textbf{扫描或错题工具} & \textbf{本作品当前实现} \\
\midrule
主要交付 & LaTeX 或 Markdown & 增强图片、PDF、文本、错题列表 & 卡片、练习、验算日志、批改、讲解与出处 \\
识别角色 & 核心能力 & 输入辅助 & 可替换入口，不宣称优于专用识别服务 \\
知识组织 & 通常不负责 & 以题或文件为主 & 以卡片为共同材料，来源与修订可查 \\
答案质量 & 不适用 & 多依赖原答案或模型输出 & 可核算题走 sympy，其他题按实际级别标注 \\
错误反馈 & 不适用 & 正误或解析展示 & 快速判分加分步批改，首个错误可定位 \\
追问依据 & 不适用 & 自由文本或联网回答 & 本地卡片召回，回答附引用，无命中不硬凑 \\
复现条件 & 依赖云服务与账号 & 依具体产品 & 缺省离线 mock，公开测试与结果文件 \\
\bottomrule
\end{tabularx}
\end{table}

上述差异不意味着本作品在所有维度更好。专用识别服务在手写公式和复杂版面上的积累更深，成熟错题产品在用户规模、同步与移动端体验上更完整。本作品的集中贡献是把识别之后的多个学习动作连成可检查的数据路径，并对最容易误导用户的答案正确性和运行状态给出证据。
